% =====================================================================
%  journal_manuscript_v2  --  LAYER-0 (spine + theorem blocks)
%  ---------------------------------------------------------------
%  Status:     Layer-1 in progress (E22-E27 ported; V6/V7 executed;
%              metric locked). The frozen v21 (scripts/
%              journal_manuscript.tex, 10,830 lines, DO NOT MODIFY)
%              remains the audit baseline.
%  Policy:     single coherent manuscript (single-paper route, secured
%              by the V5 recalibration). No diary, no version history,
%              no session references. Reads as one investigation.
%  Hierarchy:  [R-n] Result, [D-n] Definition/Design, [E-n]
%              Experiment/Evidence blocks. Every claim is tagged.
%  Sources:    Theorem B corrected per
%              download/TheoremB_Verification_and_Strengthening.md
%              (machine-verified; the (p-2)/TV/W1 defects repaired).
% =====================================================================
\documentclass[11pt,a4paper]{article}

\usepackage{amsmath,amssymb,amsthm}
\usepackage{mathtools}
\usepackage{mathrsfs}
\usepackage[margin=2.4cm]{geometry}
\usepackage{microtype}
\usepackage{booktabs}
\usepackage{tabularx}
\usepackage{enumitem}
\usepackage{xcolor}
\usepackage{graphicx}
\graphicspath{{../download/}}
\usepackage[colorlinks=true,linkcolor=blue!55!black,citecolor=blue!45!black,urlcolor=blue!55!black,bookmarks=true,bookmarksnumbered=true,unicode=true]{hyperref}

\theoremstyle{plain}
\newtheorem{theorem}{Theorem}[section]
\newtheorem{proposition}[theorem]{Proposition}
\newtheorem{lemma}[theorem]{Lemma}
\newtheorem{corollary}[theorem]{Corollary}
\theoremstyle{definition}
\newtheorem{definition}[theorem]{Definition}
\newtheorem{construction}[theorem]{Construction}
\newtheorem{assumption}[theorem]{Assumption}
\newtheorem{conjecture}[theorem]{Conjecture}
\theoremstyle{remark}
\newtheorem{remark}[theorem]{Remark}

\newcommand{\kk}{\kappa_{\!V}}
\newcommand{\kmu}{\kappa^{\mu}}
\newcommand{\HH}{\mathcal{H}}
\newcommand{\OO}{\mathcal{O}}
\newcommand{\CC}{\mathcal{C}}
\newcommand{\RR}{\mathbb{R}}
\newcommand{\TV}{\mathrm{TV}}
\newcommand{\Lip}{\operatorname{Lip}}
\newcommand{\dvol}{\,\mathrm{dvol}}
\DeclareMathOperator{\KR}{KR}

\hypersetup{
 pdftitle={A Measure-Theoretic Discrete Curvature Framework for
           Metabolic Gene Sensitivity: From Active-Set Geometry to
           Transcriptional Response},
 pdfauthor={Z.ai},
 pdfsubject={parametric FBA, discrete curvature measures, active-set
             stratification, metabolic gene sensitivity},
 pdfkeywords={discrete curvature, active set, parametric linear
              programming, flux rerouting, carbon depletion,
              transcriptional response, epistasis, holonomy}}

\title{A Measure-Theoretic Discrete Curvature Framework for Metabolic
       Gene Sensitivity: From Active-Set Geometry to Transcriptional
       Response}
\author{Z.ai}
\date{September 2026}

\begin{document}
\maketitle

\begin{abstract}
\noindent
Parametric flux balance analysis defines, for every metabolic network, a
piecewise-affine map from environmental and genetic parameters to optimal
fluxes. We show that the distributional second derivative of this map is a
matrix-valued Radon measure concentrated on the codimension-one strata of the
active-set complex, and that this measure --- not a smooth curvature ---
is the canonical carrier of the network's rerouting geometry: loop holonomy
scales linearly in loop size (slope $1.00$, versus the quadratic scaling of
any smooth map), and the second-order response mass concentrates
($100.0000\%$) on operational active-set switches. We prove a
refinement--resolution bridge (Theorem~B$'$): under mesh refinement the
atomic measure converges weakly to the smooth curvature density, the
dual-cell reconstruction converges strongly in $L^1$, and the total-variation
claims of earlier formulations are replaced by an exact counterexample
calculus; a resolution parameter $\sigma$ separates the discrete and smooth
regimes through the measured window $h \ll \sigma \ll L_{\mathrm{var}}$.
From the measure we derive a gene-level sensitivity metric $\kmu$ (measure
mass along a parametric trajectory) and show, on
\emph{E.~coli} (M3D microarrays, $n=424$ genes), that $\kmu$ predicts
carbon-depletion transcriptional response ($r = +0.395$,
$p = 2.6 \times 10^{-17}$; partial $r = +0.269$ controlling reference
level; metric-invariant $\rho = 0.99998$ against the lexicographic
definition), while the association does not propagate to the protein layer
--- consistent with a model in which cells transcribe the potential for
rerouting while buffering its translation. Double-knockout epistasis aligns
with active-set structure ($\rho_S = 0.865$), and genotype loops close with
phenotypic memory ($66\%$ non-reverting). The framework unifies three
previously distinct sensitivity objects as one measure at different
resolutions.
\end{abstract}

\tableofcontents
\newpage

% =====================================================================
\section{Introduction}\label{sec:intro}
% [Layer-1 port: v21 sec:intro, rewritten to the single story]

\paragraph{The object.}
Let a metabolic network with $m$ reactions be operated by flux balance
analysis with parameters $\theta \in \Theta \subset \RR^p$ (uptake bounds,
knockdown scalings) entering the constraints affinely. The optimal flux map
$v : \Theta \to \RR^m$ is continuous and piecewise affine (parametric
linear programming), with breakpoints on the active-set strata of the
constraint complex. Its distributional second derivative
$\mu = D^2 v$ is therefore not a function but a \emph{measure} --- the
discrete curvature carrier of the network's rerouting geometry.

\paragraph{The claim structure.}
This manuscript establishes, in order:
\begin{enumerate}[label=(\roman*)]
\item \textbf{[R1]} the atom structure and weak convergence of $\mu_h$
      under refinement (Theorem~B$'$, \S\ref{sec:theoremB}), with the
      exact counterexample calculus that separates what is provable from
      what is false in earlier formulations;
\item \textbf{[R2]} the two-layer separation: the codimension-one Hessian
      layer is the unbiased curvature estimator; the codimension-two
      angle-defect layer converges to the Gauss-map area density with an
      exact sec-law bias (Proposition~\ref{prop:seclaw});
\item \textbf{[R3]} the regime dichotomy: piecewise-affine maps with fixed
      kink complexes produce $O(\varepsilon)$ loop holonomy (measured slope
      $1.00$) where smooth maps produce $O(\varepsilon^2)$;
\item \textbf{[R4]} the value function $\Phi = c_{\mathrm{bio}}^\top v^*$
      is the canonical tie-break-free carrier, with shadow-price jumps
      measurable to $6$--$7$ digits against Danskin envelopes;
\item \textbf{[E-series]} the measure-derived metric $\kmu$ predicts
      carbon-depletion transcriptional response ($r = +0.395$),
      generalizes to oxygen-limitation and carbon-switch trajectories
      with matched responses ($r = +0.32$, $+0.22$; E-V7), and is
      buffered at the protein layer (\S\ref{sec:empirical}).
\end{enumerate}

\paragraph{Positioning.}
Kochanowski et al.\ demonstrated that global metabolic coordination under
catabolic limitation is largely passive at the protein/metabolite level,
driven by Crp and local metabolite concentrations. Our work asks a
different question: whether a gene-specific, network-derived rerouting
necessity predicts transcriptional response. We find that it does, and that
this transcript-level association does not propagate to the protein layer
--- consistent with a model in which cells transcribe the potential for
rerouting while buffering its translation.

% =====================================================================
\section{Notation protocol and the R/D/E hierarchy}\label{sec:notation}

\begin{itemize}
\item \textbf{[R-$n$]} blocks are proved results (theorems, propositions,
      lemmas) or machine-verified claims with reproducible artifacts.
\item \textbf{[D-$n$]} blocks are definitions, constructions, and design
      decisions (metric definitions, engine specifications).
\item \textbf{[E-$n$]} blocks are executed experiments with deposited
      artifacts (data files, figures, statistics).
\item Conjectures are labelled as such; plausible-but-unproven claims are
      demoted, never promoted.
\end{itemize}

% =====================================================================
\section{The discrete curvature measure of parametric FBA}
\label{sec:measure}

% [Layer-1 port: v21 active-set / stratification definitions]

\subsection{[D1] The measure}

\begin{definition}[Atomic curvature measure]\label{def:mu}
Let $u : \Omega \to \RR^m$ be continuous and piecewise affine on a
polyhedral complex $\CC$ covering $\Omega$. For each codimension-one facet
$F$ of $\CC$ shared by cells $\tau^-$, $\tau^+$, let
$[\nabla u]_F = \nabla u|_{\tau^+} - \nabla u|_{\tau^-}$, let $n_F$ be the
unit conormal oriented $\tau^- \to \tau^+$, and let $|F|$ be the facet
measure. The \emph{atomic curvature measure} is the matrix-valued Radon
measure
\[
\mu \;=\; D^2 u \;=\; \sum_{F} A_F \, \delta_F,
\qquad
A_F \;=\; [\nabla u]_F \otimes n_F \; |F| \;\in\; \RR^{m \times p \times p}.
\]
\end{definition}

\begin{remark}[Support]\label{rem:support}
$\mu$ is concentrated on the $(p-1)$-skeleton. The $(p-2)$-skeleton carries
a different measure (the angle-defect layer,
Proposition~\ref{prop:seclaw}); the two must not be conflated. In $p=1$ the
atoms are the slope jumps at breakpoints.
\end{remark}

\subsection{[D2] The measure-theoretic sensitivity metric}

\begin{definition}[$\kmu$]\label{def:kmu}
For a $C^1$ parameter trajectory $\theta(t)$, $t \in [0,T]$, and its
piecewise-affine optimal-flux response $v(\theta(t))$, let
$\{t_e\}_e$ be the event times at which the trajectory crosses active-set
strata, with atoms $A_e = A_{F(e)}$ evaluated on the crossed facet. Define
\[
\kmu(g) \;=\; \frac{1}{T}\sum_{e} \bigl\| A_e \bigr\| \, w_e(g),
\]
where $\|{\cdot}\|$ is the entrywise $1$-norm on
$\RR^{m \times p \times p}$ and $w_e(g) \in \{0,1\}$ selects the events
whose crossed stratum changes the activity of reactions assigned to gene
$g$ (max over its reactions, per the GPR association).
\end{definition}

\begin{remark}[$\kmu$ is the safe regime of the reconstruction]
$\kmu$ integrates measure mass along a one-dimensional trajectory --- the
telescoping regime in which total variation and mass are automatically
consistent (R2 remarks in \S\ref{sec:theoremB}); its resolution
independence (measured mass $288.77$ stable under $4\times$/$8\times$
trajectory refinement) is exactly the $L^1$-reconstruction regime of
Theorem~B$'$(iv). Its per-reaction masses are the components of the
flux-layer crease measure whose objective contraction is $D^2\Phi$
(Theorem~\ref{thm:coupling}): the metric is the primal layer of the
canonical object.
\end{remark}

\begin{remark}[Locked status, tie-breaking, and layer assignment]
\label{rem:lock}
The definition of $\kmu$ is final for the empirical claims of this
manuscript. (i) The carrier is the flux-layer signed crease measure
$D^2 v^*$ (Definition~\ref{def:mu} applied to $u = v^*$) of the
\emph{three-stage lexicographic} optimal flux (\S\ref{sec:methods});
the lexicographic tie-break (biomass, then $\ell_1$ flux, then seeded
deterministic weights) is part of the metric's definition, not an
implementation detail --- the flux layer is selection-dependent by
construction, and the tie-break is declared rather than eliminated.
(ii) The per-reaction mass is the unsigned total variation
$\frac{1}{T}\sum_e \|A_e\|_1$ along the parameter trajectory.
(iii) The value layer $D^2\Phi$ enters only as a \emph{derived
diagnostic} (Theorem~\ref{thm:coupling}:
$D^2\Phi = \sum_r c_r\, D^2 v^*_r$), never as the predictor: no claim
of this manuscript attributes transcriptional response to Alexandrov
curvature of the value function. (iv) The layer correlated with
transcriptomics is the codimension-one Hessian/crease layer of the
flux map --- not the Monge--Amp\`ere atom layer
(Proposition~\ref{prop:maatom}), which requires concavity and serves
as the dual-face diagnostic. Robustness to the declared selection is
measured, not assumed: the within-layer rank identity
$\rho(\kmu, \kk^{\mathrm{lex}}) = 0.99998$ (E-V5), and the trajectory
substitutions of E-V7 reproduce the association under the same
declared selection on two further paths.
\end{remark}

\subsection{[D3] The value-function layer (Alexandrov)}

\begin{proposition}[Existence and tie-break-freeness]\label{prop:alex}
Let $\Phi(\theta) = \max\{ c^\top v : Sv = 0,\ \ell(\theta) \le v \le
u(\theta)\}$ with $\ell, u$ affine, feasible and bounded. Then $\Phi$ is
concave (pointwise infimum of affine dual objectives over a
$\theta$-independent dual feasible set), and $-\Phi$ carries a symmetric
matrix-valued Radon Hessian measure $D^2(-\Phi)$ (Alexandrov), independent
of the optimal-flux selection rule. Concavity fails in general under OR
(isoenzyme) GPR rules for simultaneous multi-gene capacity vectors;
single-gene axes, exchange-bound families, and AND-only subnetworks
retain it. For OR-containing rules the value function remains
continuous piecewise affine and carries the signed crease measure of
Proposition~\ref{prop:semiconvex}; only the Monge--Amp\`ere layer
requires the concave class.
\end{proposition}

\begin{proposition}[Two-layer structure of $D^2(-\Phi)$]\label{prop:twolayer}
Let $f$ be convex piecewise affine. (i) The singular part of $D^2 f$ is
supported on the codimension-one facets, with density
$[\partial_n f]\,(g^+ - g^-) \otimes n$ per unit facet measure; it has no
atoms: the mass in an $\varepsilon$-neighborhood of a codimension-two
vertex is $O(\varepsilon)$ (measured log--log slope $1.000$). (ii) The
Monge--Amp\`ere measure $\det D^2 f$ is atomic at codimension-two
vertices, each atom equal to the volume of the normal fan
$|\operatorname{conv}\{\nabla f(\tau) : \tau \ni \text{vertex}\}|$
(measured: constant $0.5$ under $\varepsilon$-shrinking).
\end{proposition}

\begin{proposition}[Dual optimal face identity]\label{prop:dualface}
At a codimension-two vertex $\theta_v$ of a parametric-LP value function,
the Monge--Amp\`ere atom of $-\Phi$ equals the area of the normal fan, and
the fan is the image of the LP dual optimal face under
$\theta \mapsto U^\top y$:
\[
\det D^2(-\Phi)(\{\theta_v\})
= \bigl|\operatorname{conv}\{\nabla C_i\}\bigr|
= \bigl|\{U^\top y : (\lambda, y) \text{ optimal dual}\}\bigr|.
\]
Machine-verified on a two-parameter max-flow LP (the phenotype phase
plane in miniature): vertex refined to machine precision by exact
triple-tie solve, $\Phi(\theta_v)$ matching the LP value to
$10^{-15}$, fan area $0.235$ vs dual-face enumeration area
$0.235000068$ ($16$ direction LPs over the dual optimal face).
\end{proposition}

\begin{remark}[Layer separation and PhPP positioning]\label{rem:layersep}
The phenotype phase planes of Edwards--Ibarra--Palsson are the chamber
complex of $\Phi$; their sector corners are the atoms of the
Monge--Amp\`ere layer. The empirical $\kmu$ of
Definition~\ref{def:kmu} is a functional of the \emph{flux-map} measure
$D^2 v$ (Definition~\ref{def:mu}), which is tie-break-sensitive: in a
controlled follower-variable LP, the value function stays affine and
identical under two lexicographic tie-breaks while the follower's
second-difference mass flips $0.0100 \leftrightarrow 10^{-13}$. The
controlled relation between the value layer and the flux layer is now
a theorem (Theorem~\ref{thm:coupling} and
Corollary~\ref{conj:valueflux}); the association results of
\S\ref{sec:empirical} stand on the flux layer alone, which is where
the evidence (E-V5, E-V6) actually lives.
\end{remark}

\subsection{[D4] The value--flux coupling}

\begin{lemma}[Lex-selection regularity]\label{lem:lex}
Let $\ell, u$ be continuous piecewise affine and let $v^*(\theta)$ be
the three-stage lexicographic optimal flux (\S\ref{sec:methods}).
Then $v^*$ is well defined, continuous, and piecewise affine; its
graph over each GPR cell is polyhedral in $(v,\theta)$-space (each
lexicographic stage adjoins the graph of a piecewise-affine value
function, which is polyhedral, and an injective polyhedral projection
is piecewise affine). Measured: path segment residuals $\le 8 \times
10^{-14}$ (E-M1); $\Phi$ piecewise affine at $4.2 \times 10^{-13}$
(E-V1).
\end{lemma}

\begin{theorem}[Value--flux crease coupling]\label{thm:coupling}
$\Phi = c^\top v^*$ pointwise; hence, as symmetric matrix-valued
signed Radon measures on the parameter domain,
\[
D^2\Phi \;=\; \sum_{r} c_r\, D^2 v^*_r .
\]
Along any piecewise-affine trajectory, at every crease time $t_k$:
$\Delta\Phi'(t_k) = c^\top \Delta v'(t_k)$, and
$|\Delta\Phi'| \le \|c\|_\infty \,\|\Delta v'\|_1$.
\begin{enumerate}[label=(\roman*),leftmargin=2em]
\item \textbf{(Visibility dichotomy.)} An event with unique optima
      on both sides is a transversal optimal-vertex switch and is
      necessarily objective-moving; objective-invisible events
      ($c^\top \Delta v' = 0$, $\Delta v' \neq 0$) are exactly the
      degenerate reroutings inside $\ge 1$-dimensional optimal faces
      (mask-type events). The tie-break sensitivity of the flux layer
      is concentrated in the invisible events; the $c$-contraction is
      tie-break-free because it annihilates the degenerate layer.
\item \textbf{(Sparse-objective corollary.)} If $c =
      \gamma e_{\mathrm{bio}}$, then $D^2\Phi = \gamma\, D^2
      v^*_{\mathrm{bio}}$: the value layer is the crease measure of
      the single biomass component, and per-gene attribution from the
      value layer is structurally impossible.
\end{enumerate}
Machine-verified: identity error $0.0$ at all $12$ events and all
$4$ clusters of the iML1515 cut ($\Phi = c^\top v^*$ to $10^{-9}$);
$11$ of $12$ events objective-invisible, with flux $\ell_1$ slope
jumps up to $1.6 \times 10^{4}$ against $|c^\top \Delta v'| \le
1.5 \times 10^{-7}$; $150$ random dense-objective LPs to $1.1 \times
10^{-13}$ ($5$ events, all objective-moving --- the dichotomy); the
degenerate follower family $60/60$ with invisible kinks at identity
error $2.0 \times 10^{-13}$; two-dimensional mixed second
differences at crease-straddling boxes to $6.4 \times 10^{-16}$.
\end{theorem}

\begin{proposition}[Semiconvexity collapse; signed layer for mixed
GPRs]\label{prop:semiconvex}
A continuous piecewise-affine function is semiconvex (resp.\
semiconcave) iff it is convex (resp.\ concave) --- semiconvexity is
not an available weakening for LP value functions (measured: the
required constant blows up like $1/(2h)$; $\lambda \cdot h_{\max}
\equiv 0.500$ across $\lambda = 1 \dots 10^{12}$). Nevertheless
every continuous piecewise-affine $\Phi$, for any GPR structure,
carries a finite \emph{signed} symmetric Radon Hessian measure on
its crease skeleton with per-crease sign flags: for the OR-plus-cap
value function $\min(\max(\theta_1,\theta_2), K)$, the diagonal
crease density is positive semidefinite (eigenvalues
$\{0,\sqrt2\}$) and the cap crease negative semidefinite
($\{-1,0\}$). Concavity is required only by the
Monge--Amp\`ere layer and the dual-face identity; the signed Hessian
layer applies to all genes. Under multiplicative activity scaling
with $\ell \le 0 \le u$, $\Phi = \Psi(a(\theta))$ with $\Psi$
concave nondecreasing: AND-only (min-tree) GPRs retain concavity,
top-level OR nodes can break it, and disaggregating isozymes into
separate reactions restores affine bounds while replacing
max-semantics by sum-semantics (capacity $1$ vs $2$ at full
expression) --- a different biological model, not a regularization.
\end{proposition}

\begin{proposition}[Monge--Amp\`ere atoms are determinants, not
products]\label{prop:maatom}
For convex piecewise-affine $f$ on $\RR^2$ at a generic vertex with
incident gradients $n_1, n_2, n_3$ and codimension-one gradient jumps
$j_i$: the atom equals
$\tfrac12|\det(j_1, j_2)| =
|\operatorname{conv}\{n_i\}| =
\tfrac12 |j_1| |j_2| \sin\angle$; the product $|j_1||j_2|$
overestimates by $2/\sin\angle \ge 2$, with equality iff the jumps
are orthogonal (measured on $400$ random max-of-affine vertices:
determinant $=$ fan area to $8.9 \times 10^{-16}$; product/atom
median $3.296$, minimum exactly $2.000$; the $2/\sin\angle$ law to
$2.5 \times 10^{-12}$; in $p$ dimensions, atom $=
(1/p!)|\det(j_1,\dots,j_p)|$ for simplex fans).
\end{proposition}

\begin{remark}[Sign and normalization conventions]\label{rem:conventions}
Definition~\ref{def:mu} applied to $u = -\Phi$ is the \emph{signed
crease measure} of the value function: per facet $F$, the density
$A_F = [\nabla(-\Phi)]_F \otimes n_F\,|F|$, a symmetric matrix whose
eigenvalue signs flag the crease's convexity/concavity
(Proposition~\ref{prop:semiconvex}). This is the explicit extension
of Alexandrov's setting to arbitrary GPR structure: within the
concave class $-\Phi$ is convex, every $A_F \succeq 0$, and the
measure coincides with the Alexandrov Hessian measure; outside it,
the same definition produces the signed layer. All masses in this
manuscript are unsigned total variations (entrywise $1$-norm for
matrix atoms, trajectory sums for $\kmu$); Monge--Amp\`ere atoms are
unsigned $\tfrac12|\det(j_1, j_2)|$ equal to the normal-fan volume,
normalized to $(1/p!)|\det(j_1,\dots,j_p)|$ for simplex fans in
dimension $p$ (Proposition~\ref{prop:maatom}). Signed determinants
appear only in the concave class, where they are nonnegative.
\end{remark}

\begin{remark}[What this closes]\label{rem:closes}
Theorem~\ref{thm:coupling} is the exact form of the value--flux
relation previously left open: the empirical $\kmu$ is the
per-reaction total variation of the flux-layer crease measure whose
objective contraction is $D^2\Phi$ --- the metric is the primal
layer of the canonical object, justified by identity rather than by
rank agreement.
\end{remark}

% =====================================================================
\section{The refinement--resolution bridge (Theorem B$'$)}
\label{sec:theoremB}

Throughout this section $\Omega \subset \RR^p$ is compact, $p \ge 1$,
$u \in C^2(\Omega; \RR^m)$, $\{\mathcal{T}_h\}$ is a shape-regular family
of conforming triangulations with mesh size
$h = \max_\tau \operatorname{diam}(\tau) \to 0$, and $u_h$ is the
continuous piecewise-affine interpolant of $u$ on $\mathcal{T}_h$.

\begin{theorem}[Refinement--resolution bridge, corrected form]
\label{thm:Bprime}
Let $\mu_h = D^2 u_h$ (Definition~\ref{def:mu}) and
$\mu = D^2u \dvol$. Then:
\begin{enumerate}[label=(\roman*),leftmargin=2em]
\item \textbf{(Atoms / support.)} $\mu_h$ is a finite matrix-valued Radon
      measure concentrated on the $(p-1)$-skeleton of $\mathcal{T}_h$ with
      atoms $A_F = [\nabla u_h]_F \otimes n_F |F|$; there are no atoms on
      lower-dimensional strata.
\item \textbf{(Weak convergence.)} $\mu_h \Rightarrow \mu$ weakly as
      matrix-valued measures: for every
      $\varphi \in C_c^\infty(\Omega)$,
      $\langle \mu_h, \varphi \rangle \to \langle \mu, \varphi \rangle$.
      Moreover $\|\mu_h - \mu\|_{\KR} \to 0$ in the
      Kantorovich--Rubinstein (flat) norm.
\item \textbf{(Total variation: the honest statement.)}
      $\sup_h \|\mu_h\|_{\TV} \le C\,|u|_{C^2}|\Omega|$ (under
      quasi-uniformity) and
      $\liminf_h \|\mu_h\|_{\TV} \ge \|\mu\|_{\TV}$; but
      $\|\mu_h\|_{\TV} \to \|\mu\|_{\TV}$ is \emph{false in general} ---
      see Remark~\ref{rem:tvcounter}. It holds in the one-signed regime
      (convex or concave data restricted to one-dimensional parameter
      families), which is the regime of the parametric-LP value function
      and of $\kmu$.
\item \textbf{(No-mass-loss, correct form: $L^1$ reconstruction.)}
      On structured complexes, the per-cell reconstructed density
      $\rho_h(\mathrm{cell}) = |\mathrm{cell}|^{-1} \sum_{F \to
      \mathrm{cell}} A_F$ converges to $D^2 u$ \emph{strongly in $L^1$}
      with rate $O(h)$ for $u \in C^3$; consequently
      $\|\rho_h\|_{L^1} \to \|\mu\|_{\TV}$. On unstructured complexes the
      naive three-facet patch is not the dual cell; the general statement
      requires the discrete-duality (dual-diamond) pairing.
\item \textbf{(Resolution window.)} For the coarse-grained family
      $\mu_{h,\sigma} = \mu_h * \varphi_\sigma$: at fixed $h$, the limit
      $\sigma \to 0$ is the \emph{atomic} measure; the smooth object
      exists only as a family member at matched resolution through the
      window $h \ll \sigma \ll L_{\mathrm{var}}$ (the measured regime dial,
      \S\ref{sec:m4}).
\end{enumerate}
\end{theorem}

\begin{proof}[Proof of (ii) (two lines)]
For $\varphi \in C_c^\infty(\Omega)$,
$\langle D^2 u_h, \varphi\rangle = \int_\Omega u_h D^2\varphi \dvol$ by the
definition of distributional derivatives (no boundary terms arise).
Since $\|u - u_h\|_\infty \le C h^2 |u|_{C^2}$ on shape-regular families,
$\int u_h D^2\varphi \to \int u D^2\varphi = \langle \mu, \varphi\rangle$.
The flat-norm statement follows by density of
$\{c\,\varphi : \Lip(\varphi)\le 1\}$ in the KR dual cone and the uniform
mass bound of (iii).
\end{proof}

\begin{remark}[Exact counterexample calculus]\label{rem:tvcounter}
On the right-triangle mesh of $[0,1]^2$ with Hessian
$H = \left(\begin{smallmatrix} a & b\\ b & c\end{smallmatrix}\right)$, the
per-cell atoms are $A_{\mathrm{diag}} = b h^2 \mathbf{1}\mathbf{1}^\top$,
$A_{\mathrm{vert}} = (a-b) h^2\, e_x \otimes e_x$,
$A_{\mathrm{horiz}} = (c-b) h^2\, e_y \otimes e_y$; their \emph{sum} is
exactly $h^2 H$ (machine-verified to $0.0$), while the entrywise total
variation per cell is $|a-b| + |c-b| + 4|b|$ against the target
$|a| + |c| + 2|b|$. Hence for $u = xy$ the ratio is exactly $3 - 1/n$, for
the strictly convex $u = 2x^2 - xy + 2y^2$ it is $1.4 - 1/n$ (measured:
$2.9922$ and $1.3922$ at $n=128$), and for generic smooth maps $3.7$--$4.4$.
Total-variation convergence fails generically --- including under
convexity --- while weak convergence holds at rate $O(h^2)$ throughout.
\end{remark}

\begin{proposition}[Two-layer separation and the sec law]
\label{prop:seclaw}
Let $K_h = \sum_v \operatorname{defect}(v)\,\delta_v$ be the vertex
angle-defect measure of the polyhedral graph
$\{(x, u_h(x))\}$. Then $K_h$ converges weakly to the \emph{Gauss-map area
density}
\[
\frac{\det D^2 u}{\left(1 + |\nabla u|^2\right)^{3/2}} \dvol
\;=\; K\,\sqrt{1+|\nabla u|^2}\,\dvol,
\qquad
K = \frac{\det D^2 u}{(1+|\nabla u|^2)^2},
\]
i.e.\ curvature per projected parameter area, \emph{not} the intrinsic
Gaussian curvature. The bias factor is exactly
$\sqrt{1+|\nabla u|^2}$ (the sec law): machine-verified to $0.0\%$ excess
across $13$ probes spanning tilts $0.11$--$2.83$ and three tilt
directions. Recovering intrinsic curvature requires the graph area
correction.
\end{proposition}

\begin{remark}[Consequences]
(i) The codimension-one Hessian layer is the unbiased estimator and the
canonical carrier for the bridge; (ii) Regge/angle-defect formulations
need the area correction of Proposition~\ref{prop:seclaw}; (iii) defects
of a \emph{fixed} complex (as measured in the $M4b$ census) are
scale-invariant $O(1)$ objects and are not contradicted by
Proposition~\ref{prop:seclaw}, which concerns refinement limits.
\end{remark}

\begin{proposition}[Regime dichotomy; M4a in closed form]
\label{prop:dichotomy}
Let $B(\varepsilon)$ be the curvature mass captured by a loop of radius
$\varepsilon$. For a $C^2$ map, $B(\varepsilon) = O(\varepsilon^2)$; for a
piecewise-affine map whose kink complex meets the loop transversally,
$B(\varepsilon) = O(\varepsilon)$. Measured slopes on the
$M4a$ design: $2.000$ (smooth) and $1.00$ (parametric FBA,
$76$-pair perturbation scan).
\end{proposition}

\subsection{[R-port] Bridge to parametric FBA}

\begin{conjecture}[mpLP refinement bridge]\label{conj:bridge}
Let $\{\Phi_n\}$ be a sequence of parametric-LP value functions (minima of
tangent-plane families; $\theta$ in the right-hand side; nested outer
approximations) converging uniformly to a $C^2$ limit $f$, with induced
complexes $\mathcal{C}_n$. Then $\mu_n = D^2\Phi_n$ converges weakly to
$D^2 f \dvol$ provided the complexes satisfy a per-cell discrete-duality
identity (structured complexes) or a stable discrete-duality
(dual-diamond) pairing. Along one-dimensional parameter cuts the
conclusion holds unconditionally by monotone telescoping of the concave
restriction.
\end{conjecture}

\begin{remark}[Status]
The one-dimensional cut case is machine-verified (independent-seed
reproduction of the refinement experiment: folded Wasserstein distance
$0.058 \to 0.0039$, mass ratio $0.939 \to 1.0009$ across
$n = 4 \to 128$). The general complex case is open; the discrete-duality
pairing stability is the identified missing assumption.
\end{remark}

\begin{remark}[Existence vs.\ resolution]
Proposition~\ref{prop:alex} supplies existence of the curvature measure
without refinement sequences, smoothing, or interpolation; Theorem~B$'$
governs what a discretely sampled trajectory can \emph{resolve} of that
measure, and at what rate. The two layers are complementary: the
Alexandrov layer removes the extrinsic hypotheses, the refinement
bridge keeps the quantitative rate.
\end{remark}

\begin{corollary}[Decoupling; resolves the value--flux layer
relation]\label{conj:valueflux}
Theorem~\ref{thm:coupling} resolves the layer relation exactly: value
events are a subset of flux events, never the converse, and the flux
events invisible to the value layer are the degenerate reroutings
(mask-type events --- rerouting without growth loss, the
RC6/Kochanowski pattern). On the iML1515 cut: $12$ flux events vs $1$
value atom, $11$ objective-invisible; on the E24 trajectory (E-V6) the
value layer carries no network events at all. The association of
$\kmu$ with transcriptional response is therefore a flux-layer
property by structural necessity, not merely by empirical choice.
\end{corollary}

% =====================================================================
\section{Computational validation}\label{sec:computational}

% [Layer-1 port: M1/M3/M4 sections with figures from
%  download/m1_m3/ and download/m4/]

\subsection{[E-M1] Active-set curvature verification}
Twelve parameter sweeps (glucose, oxygen, eight gene knockdowns) on
\emph{E.~coli} iML1515 with an iJO1366 replication, solved by a
three-stage lexicographic pFBA engine (deterministic optimum; fixed seeded
tie-break weights after the discovery of pFBA vertex degeneracy).
Result: $100.0000\%$ of second-order response mass concentrates on
operational active-set switches in every sweep that crosses
critical-region boundaries (mass $0.934$--$1.0$; AUC $0.83$--$1.00$;
MWU $p \le 10^{-4}$); paths are piecewise affine to machine precision
between events (relative segment residuals $\le 8 \times 10^{-14}$);
negative controls behave as theory dictates (unused-pathway knockdown:
zero events, zero curvature; single-critical-region sweep: noise-level
$D^2 \sim 10^{-11}$).

\subsection{[E-M3] Double-knockout epistasis and path dependence}
$1{,}516$ single knockouts and $2{,}779$ double-knockout pairs across five
panels. Epistasis $\varepsilon_{ij} = \kappa_{ij} - \kappa_i - \kappa_j$:
all $40$ synthetic-lethal pairs are isozyme redundancies
(tktA/tktB, acnA/acnB, metE/metH, \dots) with pure-emergence epistasis;
$|\varepsilon_{ij}|$ aligns with active-set footprint overlap
(Spearman $\rho_S = 0.865$, $J_{\mathrm{support}} = 0.800$); sequential
($L_1$-MOMA) knockouts open nonzero commutators in $25\%$ of active pairs;
closed genotype loops fail to return to the initial state in $66\%$ of
pairs (phenotypic memory; median holonomy $\sim 110$).

\subsection{[E-M4] Regime dial}
\label{sec:m4}
The smoothing identity $D^2(v * \phi_\sigma) = (D^2 v) * \phi_\sigma$
holds to kernel self-test error $1.2 \times 10^{-6}$; measured crossover
$\varepsilon^*/\sigma \in [2.45, 4.11]$ (median $3.1$) separates the
$O(\varepsilon^2)$ smooth regime from the $O(\varepsilon)$ discrete
regime. The $M4a$ scaling slope is $1.00$ on the interacting stratum;
the $M4b$ census resolves the thin-wedge sliver structure
(net measure jump $8.9$ vs.\ nominal $\pm 1884.6$).

\subsection{[E-V1] The value function as canonical carrier}
$\Phi = c_{\mathrm{bio}}^\top v^*$ is single-valued with no tie-breaking;
it is piecewise affine at $4.2 \times 10^{-13}$ on the flux-event
partition, carries one real atom
$\Delta\Phi' = -0.006439$ at $t = 0.0358286$, and decouples from the
flux-jump hierarchy (value/flux mass ratio $1.7 \times 10^{-6}$).
Danskin: finite-difference shadow prices match envelope marginals to
$6$--$7$ digits. A controlled follower-variable LP reproduces the same
structure: the value function stays affine and identical under two
lexicographic tie-breaks while the follower's second-difference mass
flips $0.0100 \leftrightarrow 10^{-13}$ --- the value layer is
canonical (Proposition~\ref{prop:alex}), the flux-strain layer is not
(Corollary~\ref{conj:valueflux}). The coupling identity of
Theorem~\ref{thm:coupling} is machine-verified on this same cut:
$\Delta\Phi' = c^\top \Delta v'$ with error $0.0$ at all $12$
events and all $4$ clusters, $\Phi = c^\top v^*$ to $10^{-9}$, and
the $11$ invisible events carry flux $\ell_1$ slope jumps up to
$1.6 \times 10^{4}$ against $|c^\top \Delta v'| \le 1.5 \times
10^{-7}$ --- the decoupling is the objective contraction, not a loss
of network information.

% =====================================================================
\section{Empirical association and protein-layer buffering}
\label{sec:empirical}

% [Layer-1 port: E22-E27 sections]

\subsection{[E-E22] The gene panel and GPR classes}
Reaction-based sampling on the E22 physiology (iJO1366): $438$ of
$2{,}583$ reactions active ($17.0\%$), $435$ genes with at least one
active reaction. GPR complexity classes over active reactions:
single-gene $254$, isozyme-OR $96$, complex-AND $36$, mixed $18$,
no-GPR $34$ --- the classes that drive the concavity analysis of
Proposition~\ref{prop:semiconvex}. Panel construction is consistent
with the precursor metric family to the digit (per-gene values and
panel-level correlations reproduced exactly); the precursor geometric
metric $\kk$ is weak at the panel level ($r$ from $-0.063$ unmasked
to $+0.084$ masked) --- the discrepancy that motivated the
recalibration to $\kmu$ (E-V5).

\subsection{[E-E23] Multi-condition support}
Nine FBA conditions (baseline carbon decline; galactose, trehalose,
glycerol-3-P, and proline-N carbon sources; anaerobic-nitrate,
oxidative, and osmotic designs): active-reaction support $438$--$454$
per condition, union $537$ ($20.8\%$); genes with nonzero $\kk$
expand $435 \to 525$ across conditions. The endogeneity audit is
clean (feed caps non-binding except the designed glycerol-3-P
limitation; imposed-burden magnitudes flagged as not interpretable).
The rerouting metric is well posed across carbon sources and stress
designs: its support \emph{expands} with condition diversity rather
than migrating.

\subsection{[E-V5] Transcriptional response under carbon depletion}
On the deterministic lexicographic trajectory (E22 physiology, iJO1366),
$\kmu$ (Definition~\ref{def:kmu}) versus log-fold carbon-depletion
response (M3D, $n = 424$ genes):
Pearson $r = +0.395$ ($p = 2.6 \times 10^{-17}$), Spearman
$\rho_S = +0.414$, partial $r = +0.269$ controlling reference level
($p = 1.8 \times 10^{-8}$), top-decile versus bottom-decile response
$1.92$ vs $0.89$ (MWU $1.3 \times 10^{-7}$). Baseline reproduced to the
digit ($r = +0.3739$, $n = 433$); metric invariance
$\rho(\kmu, \kk^{\mathrm{lex}}) = 0.99998$ --- both metrics sample the
same flux-event structure along the same lexicographic trajectory, so
the association is an event-structure property of the flux layer, not a
metric artifact and not a value-function shadow. Cross-platform
replication with the locked metric: the carbon-switch arms replicate
on the same platform (M3D microarray: glycerol $+0.195$, acetate
$+0.166$, proline $+0.175$, switch-MAX $+0.206$; $n = 424$, all
$p < 6 \times 10^{-4}$) but dissociate across platforms --- the
PRECISE RNA-seq carbon-switch statistic (ten WT non-glucose carbon
conditions) gives $r = -0.044$ (partial $-0.088$, NS), with
individual conditions weakly positive (glycerol $+0.126$, acetate
$+0.130$, fructose $+0.099$) except galactose ($-0.088$). The
dissociation previously established with the precursor $\kk$
predictor is metric-robust: the association is specific to the
carbon-exhaustion class and the platform-matched design, reported as
an honest negative.

\subsection{[E-V6] The layer decision}
All value-layer attribution arms, on the frozen V5 protocol (iJO1366,
E22 physiology, $n = 424$): the flux layer $\kmu$ gives $r = +0.395$,
partial $r = +0.269$ ($p = 1.8 \times 10^{-8}$); the $c$-attribution
arm is supported on the biomass pseudo-reaction alone ($0$ of $433$
genes; Theorem~\ref{thm:coupling}(ii)); the shadow-price-jump arm
gives $51$ genes at $r = +0.032$ (partial $-0.013$, $p = 0.93$). The
value trajectory has \emph{zero} chamber crossings: the path stays in
one chamber of $\Phi$ throughout --- $\Phi$ is affine along it, with
$d\mu/dq_{\mathrm{glc}} = Y = 0.099544$ ($= |y_{\mathrm{glc}}|$, the
constant glucose uptake shadow price; affine intercept $-0.0124$) and
uptake shadow prices jumping by exactly zero --- all four value kinks
are anchor corners of the trajectory design, of magnitude
$Y\,\Delta(dq_{\mathrm{glc}}/dt)$, carrying no network structure. Value/flux strain mass ratio
$1.45 \times 10^{-3}$; the coupling identity holds at machine
precision (error $0.0$ at every kink). The value-gated flux strain
coincides with $\kmu$ exactly ($433/433$ genes attain their maximum
at a value-kink time), a property of the trajectory design (anchor
rerouting dominates per-gene maxima), not evidence of value-layer
content. The layer decision: the association lives on the flux layer,
with Theorem~\ref{thm:coupling} as the exact bridge to the canonical
object.

\subsection{[E-V7] Path robustness: what survives a trajectory change}
Three trajectories through the same LP (frozen engine, seed, panel,
statistics; only the path and its matched response vary): P0 glucose
decline (E22 anchors; the V6 control), P1 oxygen limitation
($q_{\mathrm{glc}} = 5$ fixed, $q_{O_2}$ linear from $22$ to $1$),
P2 acetate switch ($q_{O_2} = 22$ fixed; glucose $5 \to 0$ while
acetate uptake ramps $0 \to 10$). Kinks are classified by the
chamber test --- the parameters are the uptake bounds, so
$\nabla\Phi = (y_{\mathrm{glc}}, y_{O_2}, y_{\mathrm{ac}})$: a kink is
a \emph{design corner} iff it sits at a trajectory anchor and the
uptake shadow prices are continuous across it (jump $\le 10^{-9}$),
and a \emph{chamber crossing} iff they jump.
\begin{enumerate}[label=(\roman*),leftmargin=2em]
\item The one-chamber property is \emph{path-dependent}: P0 has zero
      chamber crossings (V6 reproduced digit-exactly), while P1 has
      three genuine crossings ($q_{O_2} \approx 9.3$, $6.3$, $2.9$;
      $|\Delta\Phi'|$ up to $0.68$; $y_{O_2}$ jumps $+0.032$ at the
      respiratory boundary) and P2 has four (two large, near the
      glucose--acetate handoff; two micro-slivers with dual jumps
      $\sim 10^{-4}$). The E24 single-chamber structure is a property
      of the glucose-decline design, not a law.
\item The layer decision survives the change: on every path the
      value-layer arms contribute nothing. The $c$-attribution arm is
      degenerate ($0$ of $433$ genes --- the sparse-objective
      corollary of Theorem~\ref{thm:coupling} is path-free); the
      chamber-crossing shadow-price arm is null ($r = -0.17$,
      $n = 68$ on P1; $r = +0.03$, $n = 71$ on P2); value-gated flux
      strain is a subset of $\kmu$ (equal or weaker on every path,
      never stronger).
\item The flux-layer association generalizes: against the matched
      oxygen response (M3D WT anaerobic vs aerobic contrast) P1 gives
      $r = +0.318$ (partial $+0.258$, $p = 6 \times 10^{-8}$,
      $n = 426$); against the matched acetate-switch response P2
      gives $r = +0.223$ (partial $+0.160$, $p = 9 \times
      10^{-4}$); and the P1/P2 predictors correlate with the E24
      carbon response at $r = +0.378$ / $+0.391$ --- the per-gene
      ranking is largely trajectory-independent.
\item Theorem~\ref{thm:coupling} holds at every kink on all three
      paths (identity error $0$; worst case $1.6 \times 10^{-15}$).
      Value/flux strain mass ratios: $1.5 \times 10^{-3}$ (P0),
      $2.4 \times 10^{-4}$ (P1), $5.5 \times 10^{-4}$ (P2).
\end{enumerate}

\subsection{[E-E25] Platform-by-class disambiguation ($2 \times 2$)}
Is the association a carbon-depletion artifact, or a platform
artifact? Four cells (two platforms $\times$ two perturbation
classes), executed with the precursor $\kk$ metric (rank-equivalent
to $\kmu$ at $\rho = 0.99998$): carbon switch on microarray (M3D
glycerol/acetate/proline vs glucose) switch-MAX $r = +0.196$; stress
on microarray (heat shock, acid shock, ciprofloxacin) MAX $r =
+0.298$; carbon starvation on RNA-seq (GSE64021 time course) MAX $r =
+0.191$; heat shock on RNA-seq MAX $r = +0.345$ ($n = 241$--$433$;
permutation $p \le 0.003$). Signed responses flip sign (stress arms
$-0.22$ to $-0.32$): the metric tracks \emph{rerouting magnitude},
not direction, on both platforms and both classes --- the association
is neither platform-specific nor carbon-class-specific.

\subsection{[E-E26] Protein abundance versus protein change}
Two dissociations at the protein layer (PaxDb integrated abundance;
GSE64021 quantitative proteomics; $n = 169$--$429$): (i) abundance
\emph{level} associates (PaxDb baseline $r = +0.334$) ---
high-rerouting genes are abundant proteins; (ii) abundance
\emph{change} does not (protein log-fold change $r = +0.008$ to
$+0.032$; transcript--protein fold-change coupling $r = +0.020$,
$n = 799$). The transcriptional association stops at translation,
consistent with buffering.

\subsection{[E-E27] Protein-layer null}
Condition-dependent quantitative proteomics ($22$ conditions, triplicates,
$n = 366$): protein-level correlation $r = -0.083$ against transcript
$+0.419$ on the same genes. The transcript-level association does not
propagate to the protein layer; buffering at translation is the consistent
model, positioning Kochanowski et al.\ as prior art for the protein-layer
null.

% =====================================================================
\section{Discussion: one measure, several resolutions}
\label{sec:discussion}

The three sensitivity objects of the earlier literature --- the geometric
$\kk$, the flux-statistical $\kappa_{\mathrm{flux}}$, and the time-course
$\kappa_{\mathrm{time}}$ --- are one measure at different resolutions: the
atomic measure of Definition~\ref{def:mu}, its coarse-grainings
$\mu * \phi_\sigma$, and its integrals along parameter trajectories
($\kmu$). The smooth member exists only through the window
$h \ll \sigma \ll L_{\mathrm{var}}$; the $\sigma \to 0$ limit is the
atomic object. Open items, demoted to their honest status:
Conjecture~\ref{conj:bridge} (general complex stability), E28 under the
$(\varepsilon, \sigma)$ design law, and E32 (event-measure stabilization
--- the $L^1$-reconstruction regime of Theorem~\ref{thm:Bprime}(iv)).
The formal identity $\kappa_{\mathrm{flux}} = \mathcal{F}[\mu]$ is closed by
Theorem~\ref{thm:coupling} (the exact objective contraction
$D^2\Phi = \sum_r c_r\, D^2 v^*_r$), which supersedes the earlier
metric-invariance argument for why the association stands on the flux
layer.

\paragraph{Limitations.}
Correlational design; single-organism scope; operational (not
combinatorial) active sets; lexicographic tie-breaking is an engine
convention (the value-function carrier is immune, and the flux layer
declares it --- Remark~\ref{rem:lock}); the one-chamber structure of
the E24 path is condition-specific (E-V7) while the layer decision
itself is path-robust; the cross-platform carbon-switch dissociation
(PRECISE, E-V5) bounds the generality of the association; categorical
and homotopy-theoretic material of the predecessor manuscript is
retained only where it carries falsifiable content, with the
remainder moved to a companion document.

\paragraph{What this paper is and is not.}
The mathematics is classical: parametric linear programming (chamber
complexes, active sets, Danskin duality), Alexandrov's theorem, and
distributional second derivatives of piecewise-affine maps;
Theorem~\ref{thm:coupling} is an elementary identity. The contribution
is the systematic \emph{application} of this machinery to flux-balance
models --- the two-layer measure structure of the value/flux pair,
the exact coupling identity, the signed extension to arbitrary GPR
structure --- together with the empirical association and its layer
decision. No new measure-theoretic theorem is claimed; the nearest
prior art (phenotype phase planes, mpLP value-function theory,
discrete Monge--Amp\`ere, Regge-type defect measures) is cited where
it applies.

% =====================================================================
\section{Methods}\label{sec:methods}
% [Layer-1 port: lex-pFBA engine spec, M3D/PRECISE provenance,
%  statistical protocols, reproducibility manifest]

\paragraph{Engine.} Three-stage lexicographic pFBA on HiGHS (lexicographic
order: biomass, then $\ell_1$ flux, then seeded deterministic weights);
vertex degeneracy of plain pFBA documented ($0.69$ flux flips between
warm-started solves) and resolved deterministically.

\paragraph{Data.} M3D microarray panels (carbon-depletion contrast;
WT anaerobic-vs-aerobic contrast; MOPS carbon-source switches);
PRECISE 2.0 evaluated (no carbon-depletion class, documented); PRECISE
WT carbon-switch conditions used for the cross-platform arm;
Schmidt et al.\ 2016 condition-dependent proteome; PaxDb integrated
abundance. SHA-256 provenance manifests deposited.

\paragraph{Reproducibility.} All experiments deposit artifacts
(npz/csv/json) with scripts; independent-seed reproductions recorded for
the refinement prototype and the stress batteries. The V6/V7 layer
decisions and the PRECISE arm re-run in under two minutes each
(\texttt{e24\_layer\_decision.py}, \texttt{v7\_path\_robustness.py},
\texttt{precise\_arm\_kappamu.py}).

% =====================================================================
\appendix
\section{Layer-1 porting map}\label{sec:portmap}

\begin{tabularx}{\textwidth}{lX}
\toprule
\textbf{v2 section} & \textbf{v21 source (frozen; do not modify)} \\
\midrule
\S\ref{sec:intro} & v21 sec 1 (rewritten to the single story) \\
\S\ref{sec:notation} & v21 sec 2 (subset: active-set notation only) \\
\S\ref{sec:measure} & v21 active-set stratification defs; new D1/D2 \\
\S\ref{sec:theoremB} & new (Theorem B$'$; replaces the smooth bridge) \\
\S\ref{sec:computational} & M1/M3/M4 report bodies \\
\S\ref{sec:empirical} & v21 E22--E27 sections (E22, E23, E25, E26
      ported; E24 $\to$ E-V5 with the PRECISE arm; E27 $\to$ E-E27;
      V6/V7 new) \\
\S\ref{sec:discussion} & v21 discussion, rewritten as resolution statement \\
\S\ref{sec:methods} & v21 methods; engine documentation \\
companion & categorical/HoTT sections of v21 (trimmed subset) \\
\bottomrule
\end{tabularx}

\end{document}
